% ethics.tex -- ethics statement for the arXiv v1 paper.
% Source: research/04-ethics-and-legal.md §3 (draft, adapted to a
% paragraph) and §5 (disclosure). Keep this file in sync with 04; adapt
% to the venue template at submission (arXiv has no ethics checklist --
% 04 §6). Do not weaken the "no third-party content" and
% "watermark-as-label" framing.

All text used in this study was generated by the authors using
open-weight models (opt-1.3b) via the MarkLLM toolkit
\citep{pan2024markllm}; no third-party content, user data, or live
vendor outputs were used, and no watermarked content produced by
commercial providers was collected or altered. The removal pipeline
evaluated here operates on text generated by the same user who owns it;
the study does not enable or endorse alteration of third-party content.
Watermarking is a label, not an access-control mechanism, so no
security control is circumvented. We disclose our findings to support
(i) realistic expectations for regulators relying on watermarking for
transparency obligations (Art.~50, Regulation (EU) 2024/1689
\citep{euaiact2024}), and (ii) the design of more robust provenance
mechanisms. Detectors are same-config open-source implementations;
commercial detectors were not probed. We do not provide live removal
services or weights tuned for evasion of specific vendor detectors
beyond what is reported. The authors' tooling is public
(\url{https://github.com/guillaumemeyer/watermarks-remover}) and was
deployed publicly before this study began; this paper formalizes
measurements of mechanisms already in production use.

% Disclosure (04 §5): the tool is already public and widely reported;
% no embargo or coordinated-disclosure obligation applies to this
% measurement -- state this if a venue asks about disclosure.
